% Generated draft: requires mathematical and source review.
\documentclass[11pt]{article}
\usepackage[margin=25mm]{geometry}
\usepackage{fontspec}
\setmainfont{Latin Modern Roman}
\usepackage{amsmath,amssymb,mathtools,unicode-math}
\setmathfont{Latin Modern Math}
\usepackage{enumitem,needspace,xurl,hyperref}
\setlength{\parindent}{0pt}
\setlength{\parskip}{5pt}
\setlength{\emergencystretch}{3em}
\newcommand{\N}{\mathbb{N}}
\newcommand{\Z}{\mathbb{Z}}
\newcommand{\Q}{\mathbb{Q}}
\newcommand{\R}{\mathbb{R}}
\newcommand{\C}{\mathbb{C}}
\newcommand{\F}{\mathbb{F}}
\newcommand{\A}{\mathbb{A}}
\newcommand{\PP}{\mathbb{P}}
\newcommand{\E}{\mathbb{E}}
\newcommand{\eps}{\varepsilon}
\newcommand{\dd}{\,\mathrm{d}}
\begin{document}
\section*{UnsolvedMath ID: 30003667}
\textbf{Source problem code:} catalogue entry
\subsection{Definitions and mathematical statement}
For terminology, let \((X,d)\) be a complete separable metric space, let \(p\in[1,\infty)\), and fix \(x_0\in X\). Define the space of Borel probability measures with finite \(p\)-th moment by
\[
\mathcal P_p(X)=\left\{\mu:\int_X d(x,x_0)^p\,\mathrm d\mu(x)<\infty\right\}.
\]
For \(\mu,\nu\in\mathcal P_p(X)\), let \(\Pi(\mu,\nu)\) denote their probability couplings and define
\[
W_p(\mu,\nu)=\left(\inf_{\pi\in\Pi(\mu,\nu)}\int_{X\times X}d(x,y)^p\,\mathrm d\pi(x,y)\right)^{1/p}.
\]
These definitions provide a possible framework for the terminology; the source does not specify \(X\), \(d\), or \(p\), and does not impose the hypotheses above as part of its question. A Wasserstein gradient-flow formulation would require a specified energy on a specified space of measures and a specified notion of gradient flow with respect to the chosen Wasserstein metric. None of these choices is supplied.

\textbf{Research question.} Which clustering or classification algorithms admit a formulation as Wasserstein gradient flows, and can one design such algorithms that remain stable as the sample size \(n\in\N\) tends to infinity?

\textbf{Unresolved qualifications.} The source does not specify the admissible class of algorithms, the data or label spaces, the sampling assumptions, the measure-valued representation of an algorithm, or whether the proposed representation concerns the exact discrete algorithm, a continuous-time interpolation, or a limiting evolution. It also does not define stability, the quantities to be controlled, the topology or metric of comparison, the mode of convergence, or whether bounds must be uniform in time or in \(n\). Thus the supplied question is a research program rather than a fully quantified mathematical assertion. A precise characterization or existence statement requires these choices; no particular choice is asserted here.
\subsection{Short English statement}
Identify clustering or classification algorithms that can be represented as Wasserstein gradient flows, and investigate designs whose stability persists as the sample size \(n\to\infty\). The source leaves the Wasserstein setting, algorithmic representation, sampling assumptions, and meaning of stability unspecified.
\subsection{Sources}
\begin{itemize}
\item Local registry: supplied UnsolvedMath source record, permanent numeric ID \(30003667\), title \emph{Stable Wasserstein Gradient Flow Algorithms for Data Analysis}, in the Oberwolfach Reports Open Problems collection. Exact source problem code: \texttt{catalogue entry}.
\item Supplied report citation: \emph{Variational Methods for Evolution (2018)}, Oberwolfach Reports, DOI \texttt{10.4171/owr/2017/54}, pp.~3185--3261. \url{https://doi.org/10.4171/owr/2017/54}. The registry gives proposed year \(2017\); no individual proposer is named. The supplied background attributes the collection to workshop open-problem discussions documented in Oberwolfach Reports. No problem-specific page is supplied.
\item Supplied literature reference: \emph{Accurate quantization by interacting particles}. \url{https://proceedings.mlr.press/v162/xu22d.html}.
\item Supplied literature reference: \emph{Variational Wasserstein clustering}. \url{https://pmc.ncbi.nlm.nih.gov/articles/PMC6377175/}.
\item Supplied literature reference: \emph{Dataset dynamics via Wasserstein gradient flows}. \url{https://arxiv.org/abs/2010.12760}.
\item Literature assessment dated exactly \texttt{2026-08-21}: the supplied curation labels the problem \texttt{partially\_solved} and states that Wasserstein variational clustering, interacting-particle optimization, and gradient-flow learning methods provide positive instances, while a general characterization and stability-preserving algorithmic theory is not known. This is supplied dated curation derived from the OWR record, not an independently verified assessment. The source expressly calls for independent verification of time-sensitive status and cited claims.
\item The registry labels the statement \texttt{corrected\_verified}, explaining that mathematical content was preserved while notation, typography, and sentence structure were repaired for publication. This is a supplied editorial qualification, not independent mathematical verification of this draft or its sources.
\end{itemize}
\end{document}
